\documentclass{article}

\usepackage[preprint]{neurips_2025}

\usepackage[utf8]{inputenc}
\usepackage[T1]{fontenc}
\usepackage{hyperref}
\usepackage{url}
\usepackage{booktabs}
\usepackage{amsfonts}
\usepackage{amsmath}
\usepackage{amssymb}
\usepackage{amsthm}
\newtheorem{proposition}{Proposition}
\usepackage{nicefrac}
\usepackage{microtype}
\usepackage{xcolor}
\usepackage{graphicx}
\usepackage{algorithm}
\usepackage{algorithmic}
\usepackage{multirow}
\usepackage{longtable}
\usepackage{subcaption}
\usepackage{listings}

% Custom commands
\newcommand{\tbd}[1]{\textcolor{red}{#1}}
\newcommand{\best}[1]{\textbf{#1}}
\newcommand{\gc}{\checkmark}
\newcommand{\rx}{\textcolor{red}{$\times$}}

\lstset{
  basicstyle=\ttfamily\small,
  breaklines=true,
  frame=single,
  language=Python,
  keywordstyle=\color{blue},
  commentstyle=\color{green!50!black},
  stringstyle=\color{orange}
}

\title{Supplementary Material:\\LSME: Latent-Steered Masked Editing\\for Discrete Diffusion Language Models}

\author{Anonymous Author(s)}

\begin{document}

\maketitle

\setcounter{section}{0}
\renewcommand{\thesection}{\Alph{section}}

% ============================================================
\section{Extended Method Details}
\label{app:method_details}
% ============================================================
% SOURCE: specs.md -- full module specifications

\subsection{LSME Sampler Pseudocode}
% SOURCE: specs.md Module 1 algorithm

Algorithm~\ref{alg:lsme_detailed} provides the full LSME sampling procedure with all implementation details, including entropy-based and suffix masking modes.

\begin{algorithm}[t]
\caption{LSME: Full Implementation with Mask Modes}
\label{alg:lsme_detailed}
\begin{algorithmic}[1]
\REQUIRE Model $f_\theta$, source tokens $\mathbf{x}_\text{src}$ (B, L), target latent $\mathbf{z}_\text{target}$ (B, D), mask ratio $r$, steps $S$, temperature $T$, mask mode $\in$ \{random, entropy, suffix\}
\STATE \textbf{// Step 1: Determine mask positions}
\IF{mask mode = random}
    \STATE $\mathbf{m} \sim \text{Bernoulli}(r)^{B \times L}$
\ELSIF{mask mode = entropy}
    \STATE $\text{logits}, \_ \leftarrow f_\theta(\mathbf{x}_\text{src}, \mathbf{z}_\text{target}, t{=}0.5, t_z{=}0)$
    \STATE $H_i \leftarrow -\sum_v p_v \log p_v$ for each position $i$ \COMMENT{entropy}
    \STATE $\mathbf{m} \leftarrow \text{top-}k(H, k{=}\lfloor r \cdot L \rfloor)$ \COMMENT{mask highest entropy positions}
\ELSIF{mask mode = suffix}
    \STATE $\mathbf{m}_i \leftarrow \mathbf{1}[i > (1-r) \cdot L]$ \COMMENT{mask last $\lfloor r \cdot L \rfloor$ positions}
\ENDIF
\STATE \textbf{// Step 2: Apply mask}
\STATE $\mathbf{x}_t \leftarrow \mathbf{x}_\text{src}$; \quad $\mathbf{x}_t[\mathbf{m}] \leftarrow [\text{MASK}]$
\STATE $t_\text{entry} \leftarrow r$
\STATE \textbf{// Step 3: Set target latent (clean)}
\STATE $\mathbf{z} \leftarrow \mathbf{z}_\text{target}$, \quad $t_z \leftarrow 0$
\STATE \textbf{// Step 4: Partial reverse diffusion}
\STATE $\{t_s\}_{s=0}^{S} \leftarrow \text{linspace}(\epsilon, t_\text{entry}, S+1)$
\FOR{$s = S-1, \ldots, 0$}
    \STATE $\text{logits}, \_ \leftarrow f_\theta(\mathbf{x}_t, \mathbf{z}, t_s, t_z)$
    \STATE $\text{logits}[\text{MASK}] \leftarrow -\infty$ \COMMENT{prevent sampling MASK token}
    \STATE $\text{logits} \leftarrow \text{logits} / T$ \COMMENT{temperature scaling}
    \STATE \textbf{// MDLM posterior}
    \STATE $\sigma_s \leftarrow -\log(1 - (1-\epsilon) \cdot t_s)$
    \STATE $\sigma_{s-1} \leftarrow -\log(1 - (1-\epsilon) \cdot t_{s-1})$
    \STATE $m_s \leftarrow 1 - e^{-\sigma_s}$; \quad $m_{s-1} \leftarrow 1 - e^{-\sigma_{s-1}}$
    \STATE $\text{probs} \leftarrow \text{softmax}(\text{logits}) \cdot (m_s - m_{s-1})$
    \STATE $\text{probs}[\text{MASK}] \leftarrow m_{s-1}$
    \STATE $\text{probs} \leftarrow \text{probs} / m_s$
    \STATE $\mathbf{x}_{s-1} \sim \text{Categorical}(\text{probs})$
    \STATE \textbf{// Copy flag: preserve unmasked positions}
    \STATE $\text{copy} \leftarrow (\mathbf{x}_t \neq [\text{MASK}])$
    \STATE $\mathbf{x}_t \leftarrow \text{copy} \cdot \mathbf{x}_t + (\neg\text{copy}) \cdot \mathbf{x}_{s-1}$
\ENDFOR
\RETURN $\mathbf{x}_t$, $\mathbf{m}$ \COMMENT{edited tokens and edit mask}
\end{algorithmic}
\end{algorithm}

\subsection{Attribute Latent Encoder}
% SOURCE: specs.md Module 2

The attribute latent encoder computes target latent vectors from pre-computed training latents:

\begin{lstlisting}
class AttributeLatentEncoder:
    def __init__(self, latent_dir, metadata_file):
        # Load pre-computed latent vectors and metadata
        self.latents = load_latents(latent_dir)
        self.metadata = load_json(metadata_file)

    def compute_attribute_centroids(self, attribute_name):
        centroids = {}
        for value in unique_values(self.metadata, attribute_name):
            indices = [i for i, m in enumerate(self.metadata)
                       if m[attribute_name] == value]
            centroids[value] = self.latents[indices].mean(0)
        return centroids

    def get_target_latent(self, attribute_name, target_value):
        centroids = self.compute_attribute_centroids(attribute_name)
        return centroids[target_value]

    def interpolate(self, z_source, z_target, alpha):
        # Spherical linear interpolation (SLERP)
        omega = arccos(cosine_sim(z_source, z_target))
        return (sin((1-alpha)*omega)/sin(omega)) * z_source \
             + (sin(alpha*omega)/sin(omega)) * z_target
\end{lstlisting}

\subsection{MMDiT Joint Attention Mechanism}
% SOURCE: specs.md + contribution.tex

The MMDiT block processes text and latent modalities through shared attention with modality-specific projections:

\begin{enumerate}
    \item \textbf{Input projection:} Text tokens are embedded via token embedding + positional encoding. The continuous latent is projected via an MLP.
    \item \textbf{Joint QKV attention:} Both modalities are concatenated along the sequence dimension for shared attention: $\mathbf{Q} = [\mathbf{Q}_\text{text}; \mathbf{Q}_\text{latent}]$, $\mathbf{K} = [\mathbf{K}_\text{text}; \mathbf{K}_\text{latent}]$, $\mathbf{V} = [\mathbf{V}_\text{text}; \mathbf{V}_\text{latent}]$.
    \item \textbf{Modality-specific FFN:} After attention, the sequence is split and each modality passes through its own feed-forward network.
    \item \textbf{Output:} Text logits over vocabulary $V$ and latent noise prediction $\hat{\boldsymbol{\epsilon}}$.
\end{enumerate}

% ============================================================
\section{Extended Experimental Details}
\label{app:exp_details}
% ============================================================
% SOURCE: experiments.tex

\subsection{Hardware Configuration}
% SOURCE: experiments.tex Table 5

\begin{table}[h]
\centering
\caption{Hardware configuration.}
\small
\begin{tabular}{ll}
\toprule
\textbf{Component} & \textbf{Specification} \\
\midrule
GPU & \tbd{NVIDIA A100 80GB} \\
CPU & \tbd{AMD EPYC 7763} \\
RAM & \tbd{256GB} \\
Storage & \tbd{2TB NVMe SSD} \\
\bottomrule
\end{tabular}
\end{table}
% TODO: Confirm hardware specifications.

\subsection{Baseline Implementation Details}
% SOURCE: experiments.tex Table 3

\begin{table}[h]
\centering
\caption{Baseline method details.}
\small
\begin{tabular}{llll}
\toprule
\textbf{Method} & \textbf{Venue} & \textbf{Key Idea} & \textbf{Code Repository} \\
\midrule
MDLM & NeurIPS 2024 & Masked diffusion LM & kuleshov-group/mdlm \\
ReMDM & arXiv 2025 & Remasking for diversity & kuleshov-group/remdm \\
LatentOps & ACL 2023 & ODE latent style transfer & guangyliu/LatentOps \\
DiffusER & ICLR 2023 & Edit-based diffusion & machelreid/diffuser \\
PLANNER & NeurIPS 2023 & Latent paragraph diffusion & apple/ml-planner \\
LD4LG & NeurIPS 2023 & Continuous latent DLM & justinlovelace/ld4lg \\
\bottomrule
\end{tabular}
\end{table}

\subsection{Evaluation Protocol}
% SOURCE: experiments.tex

For each experiment:
\begin{itemize}
  \item Edit 500 test samples per configuration.
  \item Report mean $\pm$ standard deviation across 3 random seeds.
  \item Use identical hardware for all methods.
  \item Apply the same DLM-Eval Suite metrics to all methods for fair comparison.
\end{itemize}

% ============================================================
\section{DLM-Eval Suite Implementation}
\label{app:eval_suite}
% ============================================================
% SOURCE: specs.md Module 4

\subsection{Pillar 1: Fluency}

\textbf{Perplexity (PPL).} Computed as $\text{PPL} = \exp(\text{NLL})$ using GPT-2 medium. We use a sliding window of 512 tokens with stride 256 for sequences longer than the model's context window.

\textbf{Grammar errors.} Counted using LanguageTool, averaged per sentence.

\subsection{Pillar 2: Controllability}

\textbf{Attribute accuracy.} A pre-trained DistilBERT classifier (fine-tuned on SST-2 for sentiment, custom for topic/formality) classifies edited text. Accuracy = fraction where predicted attribute matches target.

\subsection{Pillar 3: Edit Quality}

\textbf{BLEU.} Corpus-level BLEU-4 between source and edited texts using SacreBLEU.

\textbf{ROUGE-L.} Longest common subsequence F1 between source and edited texts.

\textbf{BERTScore.} Cosine similarity of BERT embeddings (bert-base-uncased) between source and edited texts. We report F1.

\textbf{Edit distance.} Normalized Levenshtein distance: $\text{Edit-Dist} = \text{Levenshtein}(\text{src}, \text{edited}) / \max(|\text{src}|, |\text{edited}|)$.

\subsection{Pillar 4: Latent Geometry}

\textbf{SSS (Semantic Smoothness Score).} See Section~4.1 of the main paper. Implementation uses all-MiniLM-L6-v2 for sentence embeddings and cosine similarity for adjacent comparison along the SLERP interpolation path.

\textbf{MTS (Monotonic Transition Score).} See Section~4.2 of the main paper. The classifier used matches the attribute being interpolated (e.g., sentiment classifier for negative $\to$ positive interpolation).

\textbf{Cluster separation.} Silhouette score and Calinski-Harabasz (variance ratio) index from scikit-learn, computed on latent vectors labeled by attribute.

\subsection{Pillar 5: Diversity}

\textbf{Self-BLEU.} For each generated text, compute BLEU against all other generated texts and average. Lower indicates more diversity.

\textbf{Distinct-$n$.} Fraction of unique $n$-grams (for $n \in \{1, 2, 3\}$) among all generated tokens.

\subsection{Pillar 6: Efficiency}

\textbf{NFE.} Number of forward evaluations of the model during generation/editing.

\textbf{TPS.} Tokens per second, measured as total output tokens divided by wall-clock time (excluding data loading).

% ============================================================
\section{LSME Pipeline Flowcharts}
\label{app:flowcharts}
% ============================================================
% SOURCE: specs.md flowcharts (adapted to LaTeX description)

\subsection{Overall LSME Pipeline}

The LSME editing pipeline proceeds as follows:

\begin{enumerate}
    \item \textbf{Input:} Original text (e.g., ``The food was terrible...'') and target attribute (e.g., ``positive'').
    \item \textbf{Tokenize:} Convert text to token IDs (B, L).
    \item \textbf{Attribute Latent Encoder:} Look up pre-computed centroid $\mathbf{z}_\text{target}$ for the target attribute.
    \item \textbf{Partial Mask:} Randomly mask $r$ fraction of positions $\to$ $\mathbf{x}_t$ with some [MASK] tokens.
    \item \textbf{Reverse Diffusion:} Run MMDiT from $t_\text{entry}$ to 0, with:
    \begin{itemize}
        \item Text input: partially masked tokens
        \item Latent input: target $\mathbf{z}_\text{target}$ (clean, $t_z = 0$)
        \item Each step: predict logits, sample from MDLM posterior, apply copy flag
    \end{itemize}
    \item \textbf{Output:} Edited text (e.g., ``The food was great!'').
\end{enumerate}

\subsection{Latent Geometry Analysis Pipeline}

\begin{enumerate}
    \item \textbf{Input:} Two attribute centroids (e.g., $\mathbf{z}_\text{negative}$, $\mathbf{z}_\text{positive}$).
    \item \textbf{SLERP Interpolation:} Compute $N$ latent points at $\alpha = [0, 0.1, \ldots, 1.0]$.
    \item \textbf{For each $\mathbf{z}_i$:} Generate text via L2T sampling; compute sentence embedding; run attribute classifier.
    \item \textbf{Compute:}
    \begin{itemize}
        \item SSS = mean cosine similarity of adjacent sentence embeddings
        \item MTS = fraction of monotonically increasing classifier scores
        \item Interpolation PPL = mean perplexity along the path
    \end{itemize}
\end{enumerate}

% ============================================================
\section{Extended Related Work}
\label{app:extended_related}
% ============================================================
% SOURCE: literature.md -- additional papers

\subsection{Additional Relevant Papers}

The following papers are relevant to LSME but were discussed in less detail in the main paper due to space constraints:

\begin{itemize}
    \item \textbf{CANDI}~\citep{candi2025}: Decouples discrete masking from continuous Gaussian noise via token identifiability. Unlike LSME, CANDI adds continuous noise at the token level (V-dimensional one-hot + noise) rather than using a compressed semantic latent channel.

    \item \textbf{EDLM}~\citep{xu2025edlm}: Uses energy functions to capture inter-token dependencies via importance sampling. Energy-based formulation naturally supports conditional generation but operates in discrete-only space without a learned latent.

    \item \textbf{STAR-LDM}~\citep{lovelace2025starlm}: Integrates latent diffusion ``thinking'' with AR generation using Sentence-T5 (768D) projected to 8 soft prompt tokens. Unlike LSME, it uses continuous diffusion for text generation rather than discrete masked diffusion.

    \item \textbf{Cosmos}~\citep{meshchaninov2025}: Autoencoder compresses text into a smooth latent space for continuous diffusion, achieving 2x faster inference while matching token-level quality. Demonstrates the value of smooth latent spaces for text generation.

    \item \textbf{BD3-LM}~\citep{arriola2025bd3lm}: Bridges AR and diffusion via block decomposition with per-block diffusion and AR across blocks. Achieves SOTA likelihoods among DLMs but does not incorporate continuous latent conditioning.

    \item \textbf{Duo}~\citep{sahoo2025duo}: Proves discrete diffusion naturally emerges from Gaussian diffusion, providing theoretical justification for coupling continuous and discrete processes in a single model.
\end{itemize}

\subsection{The Two Worlds of Text Diffusion}
% SOURCE: research_gap.tex Section 6

Our literature analysis confirms two disconnected paradigms in text diffusion:

\textbf{World 1 (Discrete Masked DLMs):} MDLM, SEDD, ReMDM, EDLM, CANDI, BD3-LM, LLaDA, Dream. These achieve strong language modeling with efficient masked cross-entropy training and scale to 8B parameters. However, they have limited controllability---only class-conditional CFG, remasking, or energy-based guidance.

\textbf{World 2 (Continuous Latent DLMs):} Diffusion-LM, LD4LG, PLANNER, STAR-LDM, LaDiR, Cosmos. These offer rich controllability via classifier gradients, CFG, reward guidance, and latent interpolation. However, they require continuous-to-discrete decoding, which introduces rounding artifacts.

LSME bridges these worlds by using discrete masked diffusion for text tokens (World~1 quality) while accessing a continuous latent channel for control (World~2 controllability).

% ============================================================
\section{Research Gap Analysis}
\label{app:gaps_detailed}
% ============================================================
% SOURCE: research_gap.tex -- detailed gap analysis

\subsection{Gaps Addressed by LSME}

\begin{table}[h]
\centering
\caption{Detailed gap resolution mapping.}
\small
\begin{tabular}{p{1.5cm}p{4.5cm}p{2.5cm}p{4cm}}
\toprule
\textbf{Gap ID} & \textbf{Description} & \textbf{LSME Component} & \textbf{Resolution} \\
\midrule
C4 & No DLM supports SDEdit-style editing with continuous latent steering & LSME algorithm & Partial mask + latent replacement + joint denoise \\
\addlinespace
CC2 & Latent-steered editing gap: no model combines masking-based editing with continuous latent steering & LSME algorithm & First SDEdit-style editing for discrete DLMs \\
\addlinespace
C3, TD5 & Latent space geometry unexplored in DLMs & SSS \& MTS & First formal metrics for DLM latent structure \\
\addlinespace
E1--E5 & No unified evaluation across controllability + quality + editing + efficiency & DLM-Eval Suite & 6-pillar framework \\
\addlinespace
CC4--5 & Evaluation completeness gap & DLM-Eval Suite & Addresses CCDD reviewer criticism directly \\
\bottomrule
\end{tabular}
\end{table}

\subsection{Gaps Not Addressed (Future Work)}

\begin{itemize}
    \item \textbf{Gap 3 (M6, CC3):} Guidance unification---combining CFG, classifier gradients, and reward guidance through a single latent channel. LSME demonstrates latent conditioning but does not explore all guidance modes.
    \item \textbf{Gap 5 (C1, TD1):} Conditional/seq2seq generation via latent conditioning. LSME focuses on editing rather than conditional generation.
    \item \textbf{Gap 2 (M5):} Joint end-to-end training of latent encoder + discrete text diffusion. LSME operates on a pre-trained model and does not address training methodology.
\end{itemize}

\bibliographystyle{plainnat}
\bibliography{references}

\end{document}
